\documentclass[12pt, letterpaper]{article}

\usepackage[utf8]{inputenc}
\usepackage[T1]{fontenc}
\usepackage[spanish]{babel}
\usepackage{geometry}
\usepackage{amsmath, amssymb, amsthm}
\usepackage{booktabs}
\usepackage{array}
\usepackage{microtype}
\usepackage{setspace}
\usepackage{parskip}
\usepackage{titlesec}
\usepackage{fancyhdr}
\usepackage{enumitem}
\usepackage{bm}
\usepackage{mathtools}
\usepackage{tcolorbox}
\usepackage{hyperref}

\tcbuselibrary{skins, breakable}

\geometry{top=2.8cm, bottom=3.0cm, left=3.2cm, right=3.2cm}
\setstretch{1.20}

\pagestyle{fancy}
\fancyhf{}
\fancyhead[C]{\small\textsc{Econometria · Gujarati --- Ejercicio 10.8}}
\fancyfoot[C]{\small\thepage}
\renewcommand{\headrulewidth}{0.4pt}
\renewcommand{\footrulewidth}{0pt}

\titleformat{\section}
  {\large\bfseries\scshape}{\thesection.}{0.6em}{}
  [\vspace{-4pt}\rule{\linewidth}{0.4pt}]
\titleformat{\subsection}{\normalsize\bfseries}{\thesubsection.}{0.5em}{}
\titleformat{\subsubsection}{\normalsize\itshape}{}{0em}{}

\newtheorem{prop}{Proposicion}
\newtheorem{lema}{Lema}
\theoremstyle{remark}
\newtheorem{obs}{Observacion}

\newtcolorbox{clave}{
  colback=white, colframe=black,
  boxrule=0.5pt, arc=3pt,
  left=8pt, right=8pt, top=6pt, bottom=6pt,
  breakable
}

\newcommand{\E}{\mathbb{E}}
\newcommand{\Var}{\operatorname{Var}}
\newcommand{\Cov}{\operatorname{Cov}}
\newcommand{\plim}{\operatorname{plim}}

\begin{document}

\begin{center}
  {\LARGE\bfseries Regresores Ortogonales:}\\[6pt]
  {\LARGE\bfseries Coeficientes, Interceptos y Varianzas}\\[6pt]
  {\LARGE\bfseries bajo $r_{23} = 0$}\\[12pt]
  {\large\itshape Derivacion completa desde las ecuaciones normales}\\[14pt]
  \rule{0.55\linewidth}{0.6pt}\\[8pt]
  {\normalsize Ejercicio~10.8 --- \textit{Econometria}, Gujarati \& Porter, 5.a ed.}\\[4pt]
  {\small\textsc{Con demostraciones matriciales, algebra escalar y comparacion de eficiencia}}
\end{center}

\vspace{10pt}

\noindent\textbf{Planteamiento.}\quad
Se tiene el modelo de regresion multiple:
\begin{equation}
  Y_i = \beta_1 + \beta_2 X_{2i} + \beta_3 X_{3i} + u_i
  \label{eq:multiple}
\end{equation}
y se supone que la correlacion entre los dos regresores es exactamente cero:
\begin{equation}
  r_{23} = \frac{\sum x_{2i} x_{3i}}{\sqrt{\sum x_{2i}^2 \sum x_{3i}^2}} = 0
  \quad \Longleftrightarrow \quad
  \sum x_{2i} x_{3i} = 0
  \label{eq:ortogonal}
\end{equation}
donde $x_{ji} = X_{ji} - \bar{X}_j$ denota la desviacion respecto a la media.
Se propone ademas estimar dos regresiones simples separadas:
\begin{align}
  Y_i &= \alpha_1 + \alpha_2 X_{2i} + u_{1i}
  \label{eq:simple1}\\
  Y_i &= \gamma_1 + \gamma_3 X_{3i} + u_{2i}
  \label{eq:simple2}
\end{align}
y se pregunta sobre la relacion entre los estimadores de los tres modelos.

\section{Inciso (a): \texorpdfstring{$\hat\alpha_2 = \hat\beta_2$}{alpha2 = beta2} y \texorpdfstring{$\hat\gamma_3 = \hat\beta_3$}{gamma3 = beta3}}

\subsection{Formulas generales del estimador MCO multiple}

Las ecuaciones normales del modelo~(\ref{eq:multiple}), escritas en desviaciones
respecto a la media (lo que elimina el intercepto del sistema), son:

\begin{align}
  \hat\beta_2 \sum x_{2i}^2 + \hat\beta_3 \sum x_{2i} x_{3i}
    &= \sum x_{2i} y_i
  \label{eq:normal1}\\
  \hat\beta_2 \sum x_{2i} x_{3i} + \hat\beta_3 \sum x_{3i}^2
    &= \sum x_{3i} y_i
  \label{eq:normal2}
\end{align}

donde $y_i = Y_i - \bar Y$. Resolviendo por la regla de Cramer:

\begin{equation}
  \hat\beta_2
    = \frac{(\sum x_{2i} y_i)(\sum x_{3i}^2)
            - (\sum x_{3i} y_i)(\sum x_{2i} x_{3i})}
           {(\sum x_{2i}^2)(\sum x_{3i}^2) - (\sum x_{2i} x_{3i})^2}
  \label{eq:beta2_general}
\end{equation}

\begin{equation}
  \hat\beta_3
    = \frac{(\sum x_{3i} y_i)(\sum x_{2i}^2)
            - (\sum x_{2i} y_i)(\sum x_{2i} x_{3i})}
           {(\sum x_{2i}^2)(\sum x_{3i}^2) - (\sum x_{2i} x_{3i})^2}
  \label{eq:beta3_general}
\end{equation}

\subsection{Simplificacion bajo ortogonalidad}

Aplicando la condicion~(\ref{eq:ortogonal}), $\sum x_{2i} x_{3i} = 0$,
en~(\ref{eq:beta2_general}):

\begin{equation}
  \hat\beta_2
    = \frac{(\sum x_{2i} y_i)(\sum x_{3i}^2) - 0}
           {(\sum x_{2i}^2)(\sum x_{3i}^2) - 0}
    = \frac{\sum x_{2i} y_i}{\sum x_{2i}^2}
  \label{eq:beta2_ortogonal}
\end{equation}

Esta es exactamente la formula del estimador MCO de la pendiente en la
regresion simple~(\ref{eq:simple1}):

\begin{equation}
  \hat\alpha_2 = \frac{\sum x_{2i} y_i}{\sum x_{2i}^2}
  \label{eq:alpha2}
\end{equation}

Por tanto:
\begin{equation}
  \boxed{\hat\beta_2 = \hat\alpha_2}
  \label{eq:igual_beta2}
\end{equation}

El mismo argumento para $\hat\beta_3$ con~(\ref{eq:beta3_general}):

\begin{equation}
  \hat\beta_3
    = \frac{(\sum x_{3i} y_i)(\sum x_{2i}^2) - 0}
           {(\sum x_{2i}^2)(\sum x_{3i}^2) - 0}
    = \frac{\sum x_{3i} y_i}{\sum x_{3i}^2}
    = \hat\gamma_3
  \label{eq:beta3_ortogonal}
\end{equation}

\begin{equation}
  \boxed{\hat\beta_3 = \hat\gamma_3}
  \label{eq:igual_beta3}
\end{equation}

\subsection{Interpretacion}

Cuando $r_{23} = 0$, los regresores no comparten informacion lineal. El efecto
de $X_2$ sobre $Y$ es completamente separable del efecto de $X_3$, sin que
ninguno contamine al otro. Cada coeficiente de pendiente del modelo multiple
coincide con el que se obtendria en una regresion simple que ignora la otra
variable.

En terminos de las ecuaciones normales, el cruce $\sum x_{2i}x_{3i}$ actua
como el ``canal de interferencia'' entre los dos estimadores. Al anularse,
el sistema~(\ref{eq:normal1})--(\ref{eq:normal2}) se desacopla:

\begin{align}
  \hat\beta_2 \sum x_{2i}^2 &= \sum x_{2i} y_i
  \label{eq:desacoplado1}\\
  \hat\beta_3 \sum x_{3i}^2 &= \sum x_{3i} y_i
  \label{eq:desacoplado2}
\end{align}

Cada ecuacion se resuelve independientemente, como si la otra variable no
existiera.

\begin{clave}
  La ortogonalidad ($r_{23} = 0$) es la condicion bajo la cual omitir una
  variable del modelo \textbf{no sesga} el coeficiente de la variable
  incluida. Este es el caso ideal en diseno experimental (variables
  controladas y ortogonales por construccion). En datos observacionales,
  $r_{23} = 0$ es excepcional.
\end{clave}

\section{Inciso (b): relacion entre los interceptos}

\subsection{Formula general del intercepto}

En cualquier regresion MCO, el intercepto se obtiene evaluando la recta
de regresion en las medias (condicion de primer orden respecto a $\beta_1$):

\begin{align}
  \hat\beta_1 &= \bar Y - \hat\beta_2 \bar X_2 - \hat\beta_3 \bar X_3
  \label{eq:beta1}\\
  \hat\alpha_1 &= \bar Y - \hat\alpha_2 \bar X_2
  \label{eq:alpha1}\\
  \hat\gamma_1 &= \bar Y - \hat\gamma_3 \bar X_3
  \label{eq:gamma1}
\end{align}

\subsection{Relaciones entre los tres interceptos}

Dado que $\hat\beta_2 = \hat\alpha_2$ y $\hat\beta_3 = \hat\gamma_3$,
sustituimos~(\ref{eq:alpha1}) en~(\ref{eq:beta1}):

\begin{align}
  \hat\beta_1
  &= \bar Y - \hat\beta_2 \bar X_2 - \hat\beta_3 \bar X_3
  \notag\\
  &= (\bar Y - \hat\alpha_2 \bar X_2) - \hat\beta_3 \bar X_3
  \notag\\
  &= \hat\alpha_1 - \hat\beta_3 \bar X_3
  \label{eq:beta1_vs_alpha1}
\end{align}

Analogamente, sustituyendo~(\ref{eq:gamma1}):

\begin{align}
  \hat\beta_1
  &= (\bar Y - \hat\gamma_3 \bar X_3) - \hat\beta_2 \bar X_2
  \notag\\
  &= \hat\gamma_1 - \hat\beta_2 \bar X_2
  \label{eq:beta1_vs_gamma1}
\end{align}

Resumiendo:

\begin{equation}
  \boxed{
    \hat\beta_1
      = \hat\alpha_1 - \hat\beta_3 \bar X_3
      = \hat\gamma_1 - \hat\beta_2 \bar X_2
  }
  \label{eq:relacion_interceptos}
\end{equation}

\subsection{Cuando coinciden los interceptos}

De~(\ref{eq:beta1_vs_alpha1}) se deduce:
\begin{itemize}[leftmargin=1.6em, itemsep=4pt]
  \item $\hat\beta_1 = \hat\alpha_1$ si y solo si $\hat\beta_3 \bar X_3 = 0$,
    es decir, si $\bar X_3 = 0$ (la media de $X_3$ es cero) o si
    $\hat\beta_3 = 0$ (la variable $X_3$ no tiene efecto sobre $Y$).
  \item $\hat\beta_1 = \hat\gamma_1$ si y solo si $\hat\beta_2 \bar X_2 = 0$,
    es decir, si $\bar X_2 = 0$ o $\hat\beta_2 = 0$.
\end{itemize}

En general, ninguna de estas condiciones se cumple con datos economicos reales,
por lo que $\hat\beta_1 \neq \hat\alpha_1$ y $\hat\beta_1 \neq \hat\gamma_1$.

\subsection{Interpretacion geometrica}

El intercepto ajusta la recta para que pase por el punto de medias
$(\bar X_2, \bar X_3, \bar Y)$. Cuando el modelo simple omite $X_3$, el punto
de medias que usa es $(\bar X_2, \bar Y)$. Dado que $\bar X_3 \neq 0$ en
general, las dos rectas de regresion no tienen el mismo origen vertical.

\begin{clave}
  Los interceptos de las regresiones simples \textbf{no coinciden} con el
  intercepto del modelo multiple, excepto en casos especiales ($\bar X_j = 0$
  o $\hat\beta_j = 0$). La relacion exacta es
  $\hat\beta_1 = \hat\alpha_1 - \hat\beta_3 \bar X_3$.
\end{clave}

\section{Inciso (c): comparacion de varianzas}

\subsection{Varianza de \texorpdfstring{$\hat\beta_2$}{beta2} en el modelo multiple}

Bajo ortogonalidad, la formula general~(\ref{eq:beta2_general}) se simplifica
a~(\ref{eq:beta2_ortogonal}). La varianza del estimador en el modelo multiple es:

\begin{equation}
  \Var(\hat\beta_2)
    = \frac{\sigma^2}{\sum x_{2i}^2}
  \label{eq:var_beta2}
\end{equation}

donde $\sigma^2 = \Var(u_i)$ es la varianza del error del modelo
multiple~(\ref{eq:multiple}). Nótese que el factor VIF desaparece porque
$r_{23} = 0$: $\text{VIF}_2 = 1/(1 - r_{23}^2) = 1$.

\subsection{Varianza de \texorpdfstring{$\hat\alpha_2$}{alpha2} en la regresion simple}

En la regresion simple~(\ref{eq:simple1}), el termino de error es:

\begin{equation}
  u_{1i} = u_i + \beta_3 X_{3i}
  \label{eq:error_simple}
\end{equation}

Esto es asi porque el verdadero modelo incluye $\beta_3 X_{3i}$, pero al
omitirlo, ese termino pasa al error. La varianza del error simple es:

\begin{align}
  \sigma_{u_1}^2
    &= \Var(u_i + \beta_3 X_{3i})
  \notag\\
    &= \Var(u_i) + \beta_3^2 \Var(X_{3i}) + 2\beta_3 \Cov(u_i, X_{3i})
  \notag\\
    &= \sigma^2 + \beta_3^2 \sigma_{X_3}^2
  \label{eq:var_error_simple}
\end{align}

donde usamos $\Cov(u_i, X_{3i}) = 0$ (supuesto de exogeneidad) y
$\sigma_{X_3}^2 = \Var(X_{3i})$. Por tanto:

\begin{equation}
  \sigma_{u_1}^2 = \sigma^2 + \beta_3^2 \sigma_{X_3}^2 \geq \sigma^2
  \label{eq:desigualdad_var}
\end{equation}

con igualdad solo si $\beta_3 = 0$.

La varianza del estimador simple es:

\begin{equation}
  \Var(\hat\alpha_2)
    = \frac{\sigma_{u_1}^2}{\sum x_{2i}^2}
    = \frac{\sigma^2 + \beta_3^2 \sigma_{X_3}^2}{\sum x_{2i}^2}
  \label{eq:var_alpha2}
\end{equation}

\subsection{Comparacion directa}

Restando~(\ref{eq:var_beta2}) de~(\ref{eq:var_alpha2}):

\begin{equation}
  \Var(\hat\alpha_2) - \Var(\hat\beta_2)
    = \frac{\sigma^2 + \beta_3^2 \sigma_{X_3}^2}{\sum x_{2i}^2}
    - \frac{\sigma^2}{\sum x_{2i}^2}
    = \frac{\beta_3^2 \sigma_{X_3}^2}{\sum x_{2i}^2}
    \geq 0
  \label{eq:diferencia_var}
\end{equation}

Por tanto:

\begin{equation}
  \boxed{
    \Var(\hat\alpha_2) \geq \Var(\hat\beta_2)
  }
  \label{eq:ineq_var}
\end{equation}

con igualdad si y solo si $\beta_3 = 0$ (la variable omitida es irrelevante).

\subsection{Mismo resultado para \texorpdfstring{$\hat\gamma_3$}{gamma3} y \texorpdfstring{$\hat\beta_3$}{beta3}}

Por simetria del argumento, el error de la regresion simple~(\ref{eq:simple2})
es $u_{2i} = u_i + \beta_2 X_{2i}$, con varianza:

\begin{equation}
  \sigma_{u_2}^2 = \sigma^2 + \beta_2^2 \sigma_{X_2}^2 \geq \sigma^2
  \label{eq:var_error_simple2}
\end{equation}

y por tanto:

\begin{equation}
  \Var(\hat\gamma_3) = \frac{\sigma_{u_2}^2}{\sum x_{3i}^2} \geq \frac{\sigma^2}{\sum x_{3i}^2} = \Var(\hat\beta_3)
  \label{eq:ineq_var_gamma}
\end{equation}

\begin{clave}
  Aunque los estimadores puntuales de las pendientes son identicos
  ($\hat\alpha_2 = \hat\beta_2$ y $\hat\gamma_3 = \hat\beta_3$), las
  varianzas \textbf{no lo son}. La regresion multiple es siempre mas
  eficiente:
  \[
    \Var(\hat\alpha_2) \geq \Var(\hat\beta_2), \qquad
    \Var(\hat\gamma_3) \geq \Var(\hat\beta_3)
  \]
  La razon es que el modelo multiple ``sabe'' que $X_3$ ya esta controlada,
  mientras que el modelo simple lo ignora y su error absorbe la influencia
  de $X_3$, inflando la varianza del estimador.
\end{clave}

\section{Sintesis}

\begin{table}[ht]
\centering
\caption{Comparacion entre modelo multiple y regresiones simples bajo $r_{23} = 0$}
\label{tab:resumen}
\small
\begin{tabular}{lccc}
\toprule
\textbf{Cantidad} & \textbf{Modelo multiple} & \textbf{Regresion simple}
                  & \textbf{Igualdad?} \\
\midrule
Coeficiente de $X_2$
  & $\hat\beta_2 = \dfrac{\sum x_{2i}y_i}{\sum x_{2i}^2}$
  & $\hat\alpha_2 = \dfrac{\sum x_{2i}y_i}{\sum x_{2i}^2}$
  & \textbf{Si} \\[10pt]
Coeficiente de $X_3$
  & $\hat\beta_3 = \dfrac{\sum x_{3i}y_i}{\sum x_{3i}^2}$
  & $\hat\gamma_3 = \dfrac{\sum x_{3i}y_i}{\sum x_{3i}^2}$
  & \textbf{Si} \\[10pt]
Intercepto
  & $\hat\beta_1 = \bar Y - \hat\beta_2\bar X_2 - \hat\beta_3\bar X_3$
  & $\hat\alpha_1 = \bar Y - \hat\alpha_2\bar X_2$
  & No (en gral.) \\[10pt]
Varianza de $\hat\cdot_2$
  & $\dfrac{\sigma^2}{\sum x_{2i}^2}$
  & $\dfrac{\sigma^2 + \beta_3^2\sigma_{X_3}^2}{\sum x_{2i}^2}$
  & No: simple $\geq$ multiple \\[10pt]
Varianza de $\hat\cdot_3$
  & $\dfrac{\sigma^2}{\sum x_{3i}^2}$
  & $\dfrac{\sigma^2 + \beta_2^2\sigma_{X_2}^2}{\sum x_{3i}^2}$
  & No: simple $\geq$ multiple \\
\bottomrule
\end{tabular}
\end{table}

\appendix

\section{Demostracion matricial bajo ortogonalidad}

La matriz de productos cruzados del modelo multiple es:
\begin{equation}
  \mathbf{X'X} =
  \begin{pmatrix}
    n & \sum X_{2i} & \sum X_{3i}\\
    \sum X_{2i} & \sum X_{2i}^2 & \sum X_{2i}X_{3i}\\
    \sum X_{3i} & \sum X_{2i}X_{3i} & \sum X_{3i}^2
  \end{pmatrix}
  \label{eq:XpX}
\end{equation}

Trabajando en desviaciones respecto a la media, la submatriz relevante es:
\begin{equation}
  \mathbf{S} =
  \begin{pmatrix}
    \sum x_{2i}^2 & \sum x_{2i}x_{3i}\\
    \sum x_{2i}x_{3i} & \sum x_{3i}^2
  \end{pmatrix}
  \label{eq:S}
\end{equation}

Bajo ortogonalidad, $\sum x_{2i}x_{3i} = 0$ y la matriz se vuelve diagonal:
\begin{equation}
  \mathbf{S} =
  \begin{pmatrix}
    \sum x_{2i}^2 & 0\\
    0 & \sum x_{3i}^2
  \end{pmatrix}
  \label{eq:S_diag}
\end{equation}

Su inversa es trivial:
\begin{equation}
  \mathbf{S}^{-1} =
  \begin{pmatrix}
    1/\sum x_{2i}^2 & 0\\
    0 & 1/\sum x_{3i}^2
  \end{pmatrix}
  \label{eq:Sinv}
\end{equation}

El estimador MCO es $\hat{\boldsymbol\beta} = \mathbf{S}^{-1}\mathbf{S}_{y}$
donde $\mathbf{S}_y = (\sum x_{2i}y_i, \sum x_{3i}y_i)'$:

\begin{equation}
  \begin{pmatrix}\hat\beta_2\\\hat\beta_3\end{pmatrix}
  =
  \begin{pmatrix}
    1/\sum x_{2i}^2 & 0\\
    0 & 1/\sum x_{3i}^2
  \end{pmatrix}
  \begin{pmatrix}\sum x_{2i}y_i\\\sum x_{3i}y_i\end{pmatrix}
  =
  \begin{pmatrix}
    \sum x_{2i}y_i/\sum x_{2i}^2\\
    \sum x_{3i}y_i/\sum x_{3i}^2
  \end{pmatrix}
  \label{eq:betahat_diag}
\end{equation}

Cada componente es exactamente el estimador de la regresion simple correspondiente.
La matriz de varianzas-covarianzas es:

\begin{equation}
  \Var(\hat{\boldsymbol\beta}) = \sigma^2 \mathbf{S}^{-1} =
  \begin{pmatrix}
    \sigma^2/\sum x_{2i}^2 & 0\\
    0 & \sigma^2/\sum x_{3i}^2
  \end{pmatrix}
  \label{eq:varmat}
\end{equation}

Los estimadores de las dos pendientes son no correlacionados entre si
($\Cov(\hat\beta_2, \hat\beta_3) = 0$), propiedad que solo se da bajo
ortogonalidad. Esto confirma que los resultados del modelo multiple y de
las dos regresiones simples coinciden \emph{en los coeficientes de
pendiente} pero difieren en las varianzas porque el $\sigma^2$ del modelo
multiple es menor que las varianzas de los errores de las regresiones simples.

\section{Magnitud del exceso de varianza en la regresion simple}

El exceso de varianza relativo del estimador simple respecto al multiple es:

\begin{equation}
  \frac{\Var(\hat\alpha_2) - \Var(\hat\beta_2)}{\Var(\hat\beta_2)}
    = \frac{\beta_3^2 \sigma_{X_3}^2}{\sigma^2}
    = \beta_3^2 \cdot \text{SNR}_{X_3}
  \label{eq:exceso_relativo}
\end{equation}

donde $\text{SNR}_{X_3} = \sigma_{X_3}^2/\sigma^2$ es la razon
senal-ruido de $X_3$. El exceso de varianza es mayor cuando:
\begin{enumerate}[leftmargin=1.8em, itemsep=4pt]
  \item $|\beta_3|$ es grande (efecto real de la variable omitida sobre $Y$).
  \item $\sigma_{X_3}^2$ es grande (mucha variacion en la variable omitida).
  \item $\sigma^2$ es pequeno (modelo con poco ruido, donde la omision
    importa relativamente mas).
\end{enumerate}

\vspace{14pt}
\noindent\rule{\linewidth}{0.4pt}

\noindent{\small\textit{Nota.} Los resultados del inciso~(a) generalizan al
caso de $k$ regresores: si todos los pares de regresores son ortogonales
(es decir, $\mathbf{X'X}$ es diagonal en desviaciones), entonces cada
coeficiente MCO del modelo multiple coincide con el de la regresion simple
de $Y$ sobre ese unico regresor. Este es el principio que subyace al diseno
de experimentos con bloques completos aleatorizados y factores ortogonales.}

\end{document}
